\chapter{Características generales de los detectores}

\section{Clasificación}
\begin{clave}[Tipos de detectores]
\begin{itemize}
\item \textbf{Activos} (lectura electrónica) / \textbf{pasivos} (sin electrónica: placa fotográfica, película radiocrómica, TLD\dots).
\item \textbf{Directos}: la señal sale directamente de la ionización del medio (gas, semiconductor).
\textbf{Indirectos}: la radiación se convierte en luz, que luego se convierte en señal eléctrica (centelleador + fotosensor).
\item \textbf{Integradores} (suman el efecto de muchas interacciones) / \textbf{contadores} (cada partícula individual).
\item Según la información: \textbf{contadores} (número de partículas, p.ej.\ Geiger--Müller),
\textbf{espectrómetros} (distribución en energía, p.ej.\ NaI, HPGe), \textbf{dosímetros} (energía neta depositada).
\item Por material: \textbf{gas, semiconductor, centelleador} (no excluyentes).
\end{itemize}
\end{clave}

\paragraph{Formación de la señal.} Si se deposita suficiente energía se crean portadores:
\[ N=\frac{\Delta E}{W} \]
donde $W$ es la energía media para crear un portador (par $e$--ion, par $e$--$h$, fotón\dots).
\begin{dato}[Escalas de energía por portador]
Gases $\approx30$ eV; semiconductores 3--10 eV (gap 1--5 eV; Si 3.6 eV, Ge 2.96 eV); centelleadores 20--500 eV por fotoelectrón ($\sim10$ eV para estados ópticos);
fonones $\sim$meV. Duración típica de los pulsos: 1 ns -- 10 $\mu$s.
\end{dato}
Cuanto menor es $W$, más portadores y mejor resolución en energía (menor fluctuación estadística relativa).

\paragraph{Materiales según bandas.} Conductores: sin gap, enorme corriente de fondo (inútiles).
Dieléctricos (diamante): gap grande, poca corriente de fuga pero señal pequeña.
\textbf{Semiconductores} (Si, Ge): gap moderado, permiten aplicar campos grandes con corriente de fuga
pequeña; requieren gran pureza.

\section{Modos de operación}
\begin{clave}[Pulso vs corriente]
\begin{itemize}
\item \textbf{Modo pulso}: cada interacción se procesa individualmente. Permite espectroscopía (amplitud $\propto$ energía) y timing.
Problema: \textbf{tiempo muerto}.
\item \textbf{Modo corriente}: se promedian las señales; se pierde la información individual. La corriente
es proporcional a la tasa de dosis. Se usa a tasas muy altas (dosimetría, monitores de haz).
\end{itemize}
\end{clave}

\subsection{Tiempo muerto}
Dos interacciones deben estar separadas un tiempo mínimo $\tau$ para dar señales distintas; si no, la
segunda se pierde o distorsiona la primera (apilamiento). Lo fija el componente más lento de la cadena:
el propio detector en un GM (decenas--cientos de $\mu$s), el ADC en sistemas con multicanal (pocos $\mu$s).

\begin{clave}[Modelos de tiempo muerto ($n$: tasa real, $m$: tasa medida)]
\begin{center}
\begin{tabular}{cc}
\textbf{No paralizable} & \textbf{Paralizable}\\[2pt]
$\displaystyle n=\frac{m}{1-m\tau}$ & $\displaystyle m=n\,e^{-n\tau}$\\[8pt]
Sucesos durante $\tau$ no se cuentan & Sucesos durante $\tau$ no se cuentan\\
ni lo prolongan; $m\to1/\tau$ & \emph{y lo prolongan}; $m\to0$ a tasas altas\\
\end{tabular}
\end{center}
A tasas bajas ($n\tau\ll1$) ambos coinciden: $m\approx n(1-n\tau)$.
En un sistema paralizable $m$ tiene un máximo $m_{max}=1/(e\tau)$ en $n=1/\tau$: una tasa medida corresponde a dos tasas reales.
\end{clave}
\begin{dato}[Ejemplos]
Paralizables: centelleadores con PMT, Si(Li)/Ge(Li) con electrónica lenta, HPGe con preamplificador RC.
No paralizables: contador Geiger, centelleadores con electrónica rápida (PET), rechazadores de apilamiento.
\end{dato}

\section{Espectro de alturas de pulso}
La mayoría de espectrómetros operan en modo pulso: la amplitud es proporcional a la energía
\emph{depositada} (que no siempre es la energía total incidente). Se representa el \textbf{espectro
diferencial} $\mathrm{d}N/\mathrm{d}H$ (cuentas por canal) frente a la amplitud/energía; el número de pulsos entre
$H_1$ y $H_2$ es el área bajo la curva.

\section{Eficiencia}
\begin{clave}[Eficiencia]
\[
\varepsilon_{abs}=\frac{\text{n.º detectadas}}{\text{n.º emitidas}}
=\underbrace{\frac{\text{n.º que llegan}}{\text{n.º emitidas}}}_{\varepsilon_{geo}=\Omega/4\pi}\times
\underbrace{\frac{\text{n.º detectadas}}{\text{n.º que llegan}}}_{\varepsilon_{int}}
\]
\begin{itemize}
\item \textbf{Geométrica}: ángulo sólido. Disco de radio $a$ a distancia $d$ en el eje:
\[ \Omega=2\pi\left(1-\frac{d}{\sqrt{d^2+a^2}}\right),\qquad \Omega\xrightarrow{d\gg a}\frac{\pi a^2}{d^2}\ \text{(inverso del cuadrado)} \]
\item \textbf{Intrínseca}: depende de $E$, $Z$, $\rho$ y espesor. Haz paralelo monoenergético: $\varepsilon_{int}=1-e^{-\mu x}$.
\item Para $\gamma$ se usa la \textbf{eficiencia de fotopico}: área neta del fotopico / $\gamma$ emitidos. Se mide con
fuentes calibradas a varias energías y se \textbf{interpola} con un ajuste (evitar extrapolar).
\end{itemize}
\end{clave}
\begin{ejemplo}
$a=2.5$ cm, $d=10$ cm $\Rightarrow\Omega=0.188$ sr, $\varepsilon_{geo}=1.5\%$. Duplicar $d$ la reduce $\approx4$ veces.
Fuentes extensas o fuera de eje: integración numérica o Monte Carlo.
\end{ejemplo}

\section{Resolución en energía}
Capacidad de distinguir dos energías próximas. La afectan: fluctuaciones estadísticas del número de
portadores (o de fotones de centelleo), colección incompleta de carga/luz y ruido electrónico.
\begin{clave}[Resolución]
\[
\mathrm{FWHM}=2\sqrt{2\ln2}\,\sigma\approx2.355\,\sigma,\qquad R=\frac{\mathrm{FWHM}}{H_0}
\]
Límite estadístico (Poisson con factor de Fano $F$):
\[
\sigma_N=\sqrt{FN},\quad N=\frac{E}{W}\ \Rightarrow\ R_{stat}=2.355\sqrt{\frac{FW}{E}}
\]
$F\approx0.1$--0.2 en gases y semiconductores ($F<1$ porque los sucesos de ionización no son independientes);
$F\approx1$ en centelleadores. Las contribuciones independientes se suman en cuadratura:
$\mathrm{FWHM}^2_{tot}=\mathrm{FWHM}^2_{stat}+\mathrm{FWHM}^2_{ruido}+\mathrm{FWHM}^2_{colección}+\dots$
\end{clave}
Cuando $R$ es pequeña (HPGe) se da directamente la FWHM en keV (p.ej.\ $\approx1.8$ keV a 1332 keV);
en centelleadores se da en \% (NaI $\approx7\%$ a 662 keV). Dos picos se separan si distan más de una FWHM.

\section{Respuesta y resolución temporal}
\begin{itemize}
\item \textbf{Respuesta temporal}: tiempo en formar la señal. Buen ``timing'' = flanco de subida rápido.
\item \textbf{Tiempo de subida} (10\%$\to$90\%), \textbf{tiempo de caída} (90\%$\to$10\%), constante $\tau$ de la caída exponencial.
\item \textbf{Resolución temporal}: anchura de la distribución de diferencias de tiempo entre dos detectores en coincidencia
(CRT, \emph{coincidence resolving time}). Se mide con TAC (analógico) o con marcas de tiempo digitales.
\end{itemize}

\section{Detección de rayos gamma: material y tamaño}
\begin{clave}[Elección del material]
$N=N_0e^{-\mu x}$ y $\mu$ cae rápidamente con la energía. Para $\gamma$ de alta energía hacen falta
materiales de \textbf{$Z$ y densidad altos}.
\begin{itemize}
\item Si a 10 keV: domina el fotoeléctrico ($\mu\approx10^2$ cm$^{-1}$); 300 $\mu$m absorben el 95\%. Bien para X.
\item Si a 1 MeV: casi todo Compton ($\mu\approx0.1$ cm$^{-1}$); en 300 $\mu$m interacciona el 0.3\%; harían falta 3 cm para el 95\% y aun así escapa el fotón dispersado.
\item Por eso: \textbf{Ge} (Z=32) para espectroscopía de alta resolución; \textbf{NaI(Tl)} (I, Z=53), \textbf{BGO} (Bi, Z=83) para alta eficiencia.
\end{itemize}
\end{clave}

\begin{clave}[Efectos del tamaño del detector]
\begin{itemize}
\item \textbf{Muy pequeño}: sólo la primera interacción. Fotopico (fotoeléctrico), continuo Compton y, si $E_\gamma>1.022$ MeV,
\textbf{pico de doble escape} ($E_\gamma-2m_ec^2$): ambos 511 keV escapan.
\item \textbf{Intermedio}: Compton múltiples (cuentas entre borde Compton y fotopico), \textbf{pico de escape simple}
($E_\gamma-m_ec^2$) y doble, y fotopico reforzado.
\item \textbf{Muy grande} (absorción total): idealmente sólo el pico de energía total. Extendido a cascadas: \textbf{espectroscopía
de absorción total} (TAS) para medir alimentaciones $\beta$.
\end{itemize}
\end{clave}

\begin{clave}[Efectos del material circundante]
Rayos X característicos (p.ej.\ Pb, 75--85 keV) por fotoeléctrico en el blindaje; pico de 511 keV por aniquilación de pares
creados fuera; \textbf{pico de retrodispersión} ($\approx$200--250 keV): a $180^\circ$ todas las energías incidentes dan
fotones dispersados con energía similar.
\end{clave}

\section{Etapas de la detección}
\begin{enumerate}
\item Deposición de energía en el medio (para medir energía: alta densidad y $Z$; para trazas: mínima perturbación).
\item Conversión en señal eléctrica (directa o indirecta); carga $\propto$ energía.
\item Amplificación (electrónica o multiplicación interna: avalancha, dínodos). Introduce ruido.
\item \textbf{Conformado}: el pulso estrecho de corriente se transforma en un pulso más ancho (reduce ruido) con máximo
redondeado en el tiempo de pico $T_P$ (facilita medir la amplitud). Si la forma no depende de la señal, amplitud $\propto$ energía.
\item \textbf{Digitalización}: amplitud con ADC; diferencias de tiempo con TDC.
\end{enumerate}
